\section{Composants Svelte}

Cette section décrit chacun des 13 composants de l'application, regroupés par domaine fonctionnel. Pour chaque composant, on indique ses \emph{props} (données reçues du parent), les \emph{stores} utilisés, et le rendu produit.

% ----------------------------------------------------------------
\subsection{Navigation et contrôles globaux}
% ----------------------------------------------------------------

\subsubsection{WordList.svelte}

\begin{description}
  \item[Rôle] Sidebar gauche : liste tous les mots du dictionnaire, regroupés par catégorie grammaticale.
  \item[Props] Aucune.
  \item[Stores] \texttt{wordData} (lecture), \texttt{selectedWord} et \texttt{selectedCategory} (écriture au clic).
  \item[Rendu] Pour chaque catégorie non vide, un titre \texttt{<h2>} suivi d'une liste de boutons. Chaque bouton contient un \texttt{HtmlContent} avec \texttt{disableHover=true} --- les mots de la sidebar ne déclenchent pas les info-bulles.
  \item[Interaction] Un clic sur un mot met à jour les stores \texttt{selectedWord} et \texttt{selectedCategory}, ce qui déclenche l'affichage des détails dans le panneau central.
\end{description}

Les catégories sont affichées dans cet ordre fixe~:

\begin{center}
\begin{tabular}{ll}
  \toprule
  \textbf{Clé} & \textbf{Label affiché} \\
  \midrule
  \texttt{nom} & Noms \\
  \texttt{adj} & Adjectifs \\
  \texttt{pron} & Pronoms (déterminants) \\
  \texttt{verb} & Verbes \\
  \texttt{conj} & Conjonctions \\
  \texttt{part} & Particules \\
  \texttt{card} & Cardinaux \\
  \texttt{proposs} & Pronoms possessifs \\
  \texttt{proper} & Pronoms personnels \\
  \texttt{prep} & Prépositions \\
  \bottomrule
\end{tabular}
\end{center}

\subsubsection{AccentCheckbox.svelte}

\begin{description}
  \item[Rôle] Case à cocher flottante (position fixe, en bas à droite) pour activer/désactiver l'affichage des accents toniques.
  \item[Props] Aucune.
  \item[Stores] \texttt{accentEnabled} (lecture et écriture).
  \item[Rendu] Un \texttt{<input type="checkbox">} lié au store. Coché par défaut.
\end{description}

% ----------------------------------------------------------------
\subsection{Fiches détaillées}
% ----------------------------------------------------------------

\subsubsection{WordDetails.svelte}

\begin{description}
  \item[Rôle] Panneau central : affiche la fiche détaillée du mot sélectionné.
  \item[Props] Aucune.
  \item[Stores] \texttt{wordData}, \texttt{selectedWord}, \texttt{selectedCategory} (lecture).
  \item[Rendu] Utilise \texttt{\$derived} pour calculer \texttt{details} à partir des stores. Selon la catégorie, délègue l'affichage à un sous-composant~:
\end{description}

\begin{center}
\begin{tabular}{ll}
  \toprule
  \textbf{Catégorie} & \textbf{Composant} \\
  \midrule
  \texttt{nom} & \texttt{NounDetails} \\
  \texttt{adj}, \texttt{proposs}, \texttt{card}, \texttt{pron} & \texttt{AdjectiveDetails} \\
  \texttt{verb} & \texttt{VerbDetails} \\
  \texttt{conj}, \texttt{part} & \texttt{BaseDetails} \\
  \bottomrule
\end{tabular}
\end{center}

Si le mot possède des phrases d'exemple (\texttt{details.phrases}), un composant \texttt{ExamplePhrases} est affiché en dessous.

\subsubsection{NounDetails.svelte}

\begin{description}
  \item[Rôle] Tableau de déclinaison d'un nom : 6 cas $\times$ 2 nombres (singulier/pluriel).
  \item[Props] \texttt{details} --- l'objet du nom dans \texttt{wordData}.
  \item[Rendu] Pour chaque cas (\ukr{називний}, \ukr{родовий}, \ukr{давальний}, \ukr{знахідний}, \ukr{орудний}, \ukr{місцевий}), une ligne avec la forme singulière et plurielle accentuée.
\end{description}

% ============ PLACEHOLDER ============
% Capture : tableau de déclinaison d'un nom (ex: "будинок")
% montrant les 6 cas avec singulier et pluriel
\begin{figure}[ht]
  \centering
  \fbox{\parbox{0.6\textwidth}{\centering\vspace{2.5cm}\textit{[screenshot-noun-table.png]\\Tableau de déclinaison d'un nom}\vspace{2.5cm}}}
  % \includegraphics[width=0.6\textwidth]{screenshot-noun-table.png}
  \caption{Déclinaison d'un nom --- 6 cas, singulier et pluriel.}
  \label{fig:noun-table}
\end{figure}

\subsubsection{AdjectiveDetails.svelte}

\begin{description}
  \item[Rôle] Tableau de déclinaison : 6 cas $\times$ 4 genres (masculin, féminin, neutre, pluriel). Utilisé pour les adjectifs, pronoms possessifs, cardinaux et pronoms déterminants.
  \item[Props] \texttt{details}.
  \item[Rendu] Un tableau avec les cas en lignes (noms ukrainiens : \ukr{називний}, \ukr{родовий}, etc.) et les genres en colonnes (\ukr{чол. р.}, \ukr{жін. р.}, \ukr{сер. р.}, \ukr{множина}).
\end{description}

\subsubsection{VerbDetails.svelte}

\begin{description}
  \item[Rôle] Tableau de conjugaison complet d'un verbe.
  \item[Props] \texttt{details}.
  \item[Stores] \texttt{wordData} (pour le couple aspectuel).
  \item[Rendu] Structure du tableau~:
  \begin{itemize}
    \item \textbf{Infinitif} (\ukr{Інфінітив}) --- forme de base accentuée.
    \item \textbf{Par temps} (\ukr{Наказовий спосіб}, \ukr{Майбутній час}, \ukr{Теперішній час}, \ukr{Минулий час}) --- pour chaque temps présent dans les données :
    \begin{itemize}
      \item Passé : lignes par genre (\ukr{чол. р.}, \ukr{жін. р.}, \ukr{сер. р.}) avec singulier et pluriel fusionné.
      \item Autres : lignes par personne (\ukr{1 особа}, \ukr{2 особа}, \ukr{3 особа}) avec colonnes \ukr{Однина}/\ukr{Множина}.
    \end{itemize}
    \item \textbf{Forme impersonnelle} (\ukr{Безособова форма}) --- si présente.
    \item \textbf{Couple aspectuel} --- si le verbe a un champ \texttt{coupl}, affiche le verbe de l'aspect opposé.
  \end{itemize}
\end{description}

% ============ PLACEHOLDER ============
% Capture : tableau de conjugaison d'un verbe (ex: "читати")
% montrant infinitif, présent, futur, passé avec toutes les personnes
\begin{figure}[ht]
  \centering
  \fbox{\parbox{0.6\textwidth}{\centering\vspace{3cm}\textit{[screenshot-verb-conj.png]\\Tableau de conjugaison d'un verbe}\vspace{3cm}}}
  % \includegraphics[width=0.6\textwidth]{screenshot-verb-conj.png}
  \caption{Conjugaison complète d'un verbe avec tous les temps disponibles.}
  \label{fig:verb-conj}
\end{figure}

\subsubsection{BaseDetails.svelte}

\begin{description}
  \item[Rôle] Affichage minimal pour les mots invariables (conjonctions, particules).
  \item[Props] \texttt{details}.
  \item[Rendu] Titre «~Forme de base~» suivi de la forme accentuée, calculée avec \texttt{firstPair(details.base)}.
\end{description}

% ----------------------------------------------------------------
\subsection{Phrases}
% ----------------------------------------------------------------

\subsubsection{PhraseList.svelte}

\begin{description}
  \item[Rôle] Page de recherche de phrases (\texttt{/phrases}).
  \item[Props] Aucune.
  \item[Stores] \texttt{phraseData} (lecture).
  \item[État local] \texttt{searchQuery} --- texte saisi dans la barre de recherche.
  \item[Rendu] Barre de recherche flottante (en haut à droite) + liste filtrée. Chaque phrase est rendue via \texttt{HtmlContent} avec sa traduction française et ses références. Le filtrage utilise \texttt{filterPhrases()} (logique AND multi-termes sur les métadonnées \texttt{ref}).
\end{description}

% ============ PLACEHOLDER ============
% Capture : page /phrases avec une recherche active (ex: "L1")
% montrant les phrases filtrées
\begin{figure}[ht]
  \centering
  \fbox{\parbox{0.7\textwidth}{\centering\vspace{2.5cm}\textit{[screenshot-phrase-search.png]\\Recherche filtrée dans les phrases}\vspace{2.5cm}}}
  % \includegraphics[width=0.7\textwidth]{screenshot-phrase-search.png}
  \caption{Recherche de phrases avec filtre actif.}
  \label{fig:phrase-search}
\end{figure}

\subsubsection{ExamplePhrases.svelte}

\begin{description}
  \item[Rôle] Affiche les phrases d'exemple associées à un mot (sous la fiche détaillée).
  \item[Props] \texttt{phrases} --- objet dont les clés sont des identifiants de phrases.
  \item[Stores] \texttt{phraseData}, \texttt{wordData} (lecture).
  \item[État local] \texttt{expandedVerbs} --- objet indiquant quelles phrases ont les formes verbales dépliées.
  \item[Rendu] Pour chaque phrase~: contenu HTML (via \texttt{HtmlContent}), traduction, remarque optionnelle. Si la phrase a un champ \texttt{genereVerbe}, une case à cocher permet de déplier toutes les formes conjuguées (générées par \texttt{generateVerbForms()}).
\end{description}

% ----------------------------------------------------------------
\subsection{Grammaire interactive}
% ----------------------------------------------------------------

\subsubsection{HtmlContent.svelte}

C'est le composant central de l'interactivité. Il reçoit du HTML brut (contenant des spans \texttt{.ukr}) et y applique deux traitements~:

\begin{description}
  \item[Props] \texttt{html} (chaîne HTML), \texttt{disableHover} (booléen, défaut \texttt{false}).
  \item[Stores] \texttt{wordData}, \texttt{accentEnabled}, \texttt{pinnedElement}, \texttt{grammarTableData}.
  \item[Fonctionnement] Un \texttt{\$effect} réactif se déclenche quand \texttt{accentEnabled} ou \texttt{html} change~:
  \begin{enumerate}
    \item \textbf{applyAccents()} --- Parcourt tous les \texttt{.ukr} spans. Si les accents sont activés, utilise \texttt{data-info} pour retrouver la position d'accent dans \texttt{wordData}, puis applique \texttt{highlightLetter()} pour colorer la lettre accentuée en rouge (\textit{crimson}).
    \item \textbf{applyHoverHandlers()} --- \emph{Seulement si} \texttt{disableHover} est \texttt{false}. Attache 3 écouteurs d'événements à chaque span \texttt{.ukr}~:
    \begin{itemize}
      \item \texttt{mouseenter} : colore le mot selon le cas grammatical + affiche la bulle d'info + met à jour \texttt{grammarTableData}.
      \item \texttt{mouseleave} : retire la couleur, cache la bulle, efface \texttt{grammarTableData} (sauf si un élément est épinglé).
      \item \texttt{click} : épingle/désépingle l'élément via \texttt{pinnedElement}.
    \end{itemize}
  \end{enumerate}
\end{description}

\paragraph{Le mécanisme d'accent} L'attribut \texttt{data-info} de chaque span \texttt{.ukr} contient une chaîne séparée par des points-virgules, par exemple~:

\begin{ukrbox}
\texttt{data-info="будинок;nom;cas;nomi;s"}
\end{ukrbox}

Cette chaîne est parsée par \texttt{parseInfo()} pour obtenir : le lemme (\ukr{будинок}), la catégorie (\texttt{nom}), puis les tokens de navigation dans le JSON (\texttt{cas}, \texttt{nomi}, \texttt{s}). La fonction \texttt{getDataFromJson()} utilise ces tokens pour retrouver l'entrée correspondante, d'où on extrait la position d'accent.

% ============ PLACEHOLDER ============
% Capture : info-bulle apparaissant au survol d'un mot ukrainien
% dans le panneau central ou dans une phrase
\begin{figure}[ht]
  \centering
  \fbox{\parbox{0.5\textwidth}{\centering\vspace{2cm}\textit{[screenshot-hover-bubble.png]\\Info-bulle au survol}\vspace{2cm}}}
  % \includegraphics[width=0.5\textwidth]{screenshot-hover-bubble.png}
  \caption{Info-bulle au survol d'un mot ukrainien : lemme accentué + métadonnées.}
  \label{fig:hover-bubble}
\end{figure}

\subsubsection{GrammarSidebar.svelte}

\begin{description}
  \item[Rôle] Sidebar droite affichant le tableau de grammaire complet du mot survolé ou épinglé.
  \item[Props] Aucune.
  \item[Stores] \texttt{grammarTableData}, \texttt{pinnedElement} (lecture).
  \item[Rendu] Visible uniquement si \texttt{grammarTableData} n'est pas \texttt{null}. Affiche un en-tête avec le lemme accentué et les métadonnées (catégorie, genre pour les noms, aspect pour les verbes), puis délègue au composant \texttt{GrammarTable}. Quand un élément est épinglé, la bordure change de couleur (classe CSS \texttt{.pinned}).
\end{description}

% ============ PLACEHOLDER ============
% Captures : sidebar grammaire en mode hover et en mode épinglé
\begin{figure}[ht]
  \centering
  \fbox{\parbox{0.4\textwidth}{\centering\vspace{2.5cm}\textit{[screenshot-grammar-sidebar.png]\\Sidebar grammaire (hover)}\vspace{2.5cm}}}
  % \includegraphics[width=0.4\textwidth]{screenshot-grammar-sidebar.png}
  \caption{Sidebar de grammaire en mode survol.}
  \label{fig:grammar-sidebar}
\end{figure}

\begin{figure}[ht]
  \centering
  \fbox{\parbox{0.4\textwidth}{\centering\vspace{2.5cm}\textit{[screenshot-grammar-pinned.png]\\Sidebar grammaire (épinglée)}\vspace{2.5cm}}}
  % \includegraphics[width=0.4\textwidth]{screenshot-grammar-pinned.png}
  \caption{Sidebar de grammaire épinglée (après clic sur un mot).}
  \label{fig:grammar-pinned}
\end{figure}

\subsubsection{GrammarTable.svelte}

\begin{description}
  \item[Rôle] Génère le tableau HTML de grammaire selon la catégorie du mot.
  \item[Props] \texttt{data} --- objet \texttt{\{word, category, infos, wordData\}}.
  \item[Rendu] Utilise \texttt{\$derived.by()} pour calculer le HTML du tableau. Selon la catégorie~:
  \begin{itemize}
    \item \textbf{Noms} : cas $\times$ nombre, forme courante en gras.
    \item \textbf{Pronoms personnels} : cas simple, forme courante en gras.
    \item \textbf{Adjectifs/pronoms/cardinaux} : cas $\times$ genre, forme courante en gras.
    \item \textbf{Verbes} : infinitif + tous les temps (présent, futur, passé, impératif, impersonnel) + couple aspectuel.
  \end{itemize}
  La fonction interne \texttt{renderCell()} affiche les formes avec accent, séparées par «~/ ~» en cas de variantes multiples.
\end{description}

\subsubsection{UkrSpan.svelte}

\begin{description}
  \item[Rôle] Version composant-propre d'un span ukrainien interactif (alternative à l'approche DOM de \texttt{HtmlContent}).
  \item[Props] \texttt{dataInfo} (chaîne), \texttt{text} (texte à afficher).
  \item[Stores] \texttt{accentEnabled}, \texttt{pinnedElement}, \texttt{grammarTableData}, \texttt{wordData}.
  \item[Rendu] Un \texttt{<span class="ukr">} avec gestion de l'accent, survol, bulle et épinglage, implémentée de manière réactive plutôt que par manipulation DOM.
\end{description}

% ============ PLACEHOLDER ============
% Captures : accents activés vs désactivés
\begin{figure}[ht]
  \centering
  \fbox{\parbox{0.4\textwidth}{\centering\vspace{2cm}\textit{[screenshot-accent-on.png]\\Accents activés}\vspace{2cm}}}
  \hfill
  \fbox{\parbox{0.4\textwidth}{\centering\vspace{2cm}\textit{[screenshot-accent-off.png]\\Accents désactivés}\vspace{2cm}}}
  % \includegraphics[width=0.45\textwidth]{screenshot-accent-on.png}
  % \hfill
  % \includegraphics[width=0.45\textwidth]{screenshot-accent-off.png}
  \caption{Comparaison : accents activés (gauche) vs désactivés (droite).}
  \label{fig:accents}
\end{figure}
